\section{Vista General del Proyecto}

El proyecto Kit Pagos Colombia surge a partir de un problema recurrente en el desarrollo de software colombiano: integrar pasarelas de pago sigue siendo un proceso fragmentado, repetitivo y, en muchos casos, innecesariamente complejo. Plataformas como Wompi, Rapyd, Mercado Pago y Kushki ofrecen servicios similares, pero cada una maneja estructuras distintas para autenticación, respuestas HTTP, validaciones y notificaciones asíncronas. En la práctica, esto obliga a los equipos de desarrollo a implementar varias veces la misma lógica, aumentando el tiempo de integración y la probabilidad de errores.

De acuerdo con la Cámara Colombiana de Comercio Electrónico, en 2025 el país llegó a 9,3 millones de compradores digitales, y las ventas en línea pesaron un 2,4 \% dentro del comercio minorista que registra el DANE (“Informe de comercio minorista / Diciembre 2025”, s. f.). La Cámara lo describe como un crecimiento sostenible, lo que muestra la necesidad de herramientas que simplifiquen la incorporación de servicios de pago dentro de aplicaciones modernas. A pesar de esto, el ecosistema colombiano continúa dependiendo de APIs heterogéneas que exigen procesos de integración independientes para cada proveedor.

Frente a este escenario, el proyecto propone la construcción de un framework open source que permita unificar la interacción con distintas pasarelas bajo una misma interfaz de programación. La intención no es reemplazar las APIs oficiales de los proveedores, sino construir una capa intermedia que reduzca la complejidad técnica que enfrentan los desarrolladores cuando necesitan trabajar con múltiples pasarelas dentro de una misma aplicación.

El framework está compuesto por cuatro elementos. El primero es un SDK de servidor con interfaz unificada que abstrae las diferencias entre Wompi, Rapyd, Mercado Pago y Kushki. Cubre las operaciones comunes: crear pagos con tarjeta o PSE, consultar su estado, listar los bancos PSE y recibir webhooks. El segundo es un módulo de navegador que tokeniza la tarjeta en la página del comprador, para que su número no pase por el servidor del comercio. El tercero es una API de Simulación desplegada, que imita las APIs de las pasarelas para probar sin depender de los entornos sandbox externos y que también expone las operaciones del kit por REST. Finalmente, el proyecto incluye una documentación centralizada con los datos de prueba y las configuraciones de cada proveedor, para reducir el tiempo que se pierde navegando entre distintas documentaciones oficiales.

El propósito principal del proyecto es desarrollar y evaluar una solución que simplifique la integración de múltiples pasarelas de pago en el contexto colombiano. Más allá de construir software funcional, el trabajo busca analizar si una arquitectura unificada realmente reduce la complejidad de implementación frente a una integración tradicional realizada directamente sobre las APIs nativas de cada proveedor.

La propuesta se desarrolla bajo el enfoque de Design Science Research (DSR), metodología orientada a proyectos cuyo resultado principal es la construcción y validación de un artefacto tecnológico (Peffers, Tuunanen, Rothenberger \& Chatterjee, 2007). Esto implica que el framework no solamente debe funcionar correctamente, sino que además debe evaluarse mediante criterios técnicos verificables. Para ello, el proyecto compara una tienda que integra las pasarelas con el framework frente a otra que lo hace directamente contra las APIs nativas, con seis variables, entre ellas las métricas WMC, CBO y RFC del conjunto de métricas CK propuesto por Chidamber y Kemerer (1994).

El diseño del framework se apoya principalmente en la arquitectura hexagonal propuesta por Cockburn (2024), la cual permite separar el núcleo de negocio de las dependencias externas mediante puertos y adaptadores. Esta decisión arquitectónica resulta especialmente útil en un ecosistema como el de las pasarelas de pago, donde los proveedores cambian constantemente sus mecanismos internos y sus convenciones de integración. Adicionalmente, se emplean patrones de diseño del catálogo GoF como Strategy, Facade y Factory para encapsular el comportamiento específico de cada pasarela y ofrecer una interfaz homogénea hacia el desarrollador (Gamma, Helm, Johnson \& Vlissides, 1994).

El alcance del proyecto está delimitado a cuatro pasarelas utilizadas activamente en Colombia: Wompi, Rapyd (que adquirió PayU GPO), Mercado Pago y Kushki, y a dos métodos de pago, tarjeta y PSE. El framework cubre operaciones esenciales relacionadas con pagos digitales y manejo de notificaciones asíncronas, además de la simulación de escenarios controlados de prueba. Sin embargo, no contempla procesamiento financiero real, integración con sistemas contables ni soporte para proveedores distintos a los cuatro definidos. Tampoco modifica los mecanismos internos de seguridad implementados por las pasarelas, ya que el framework actúa únicamente como intermediario entre el cliente y los proveedores externos.

Como resultado del proyecto se obtienen un SDK de servidor basado en arquitectura hexagonal y su módulo de navegador, una API de Simulación desplegada y documentación pública de datos de prueba. A eso se suman cuatro proyectos prototípicos, un informe comparativo de la evaluación y el repositorio público del framework, con su documentación técnica de uso y una página de presentación. La intención es que estos entregables no se limiten únicamente al contexto académico del trabajo de grado, sino que puedan servir como base reutilizable para desarrolladores y futuros proyectos relacionados con integración de pagos en Colombia.

\subsection{Visión del Producto}

El ecosistema colombiano de pasarelas de pago está fragmentado. Wompi, Rapyd, Mercado Pago y Kushki funcionan cada una con sus propias convenciones de autenticación, estructuras de respuesta distintas y mecanismos de notificación asíncrona particulares. Para un equipo de desarrollo que necesita integrar más de una, esa fragmentación se convierte en trabajo repetitivo: hay que estudiar cuatro documentaciones, implementar cuatro lógicas de autenticación y normalizar manualmente los formatos de respuesta. El tiempo invertido crece con cada proveedor adicional y la probabilidad de errores también.

El Kit Pagos Colombia busca resolver eso. Es un framework open source que unifica la integración de estas pasarelas bajo una interfaz común, de modo que un desarrollador pueda interactuar con varias de ellas usando la misma lógica de código. Además, incluye una API de simulación que elimina la dependencia de los entornos sandbox externos durante el desarrollo, y una documentación que centraliza los datos de prueba dispersos en múltiples portales.

La visión de largo plazo es que el kit funcione como referencia técnica dentro del ecosistema fintech colombiano, reduciendo las barreras de entrada para que proyectos de menor escala puedan incorporar múltiples opciones de pago sin duplicar esfuerzo. Este trabajo de grado construye y evalúa la primera versión funcional de esa visión.

\subsection{Propósito, Alcance y Objetivos}

\subsubsection{Propósito}

Este proyecto se realiza como trabajo de grado en Ingeniería de Sistemas. Su propósito es construir y evaluar un framework que resuelva el problema de la fragmentación en el proceso de integración de pasarelas de pago en Colombia. La construcción del framework se enmarca en Design Science Research, que exige no solo que el sistema sea desarrollado correctamente, sino que sea validado en escenarios representativos del problema original y que sus resultados sean comunicados de forma rigurosa.

\subsubsection{Alcance}

El proyecto incluye el diseño, desarrollo y evaluación de cuatro componentes, los mismos que describe el SRS v2.0:

\begin{itemize}
  \item El SDK de servidor (\texttt{kit-pagos-colombia}): una biblioteca TypeScript que oculta las diferencias entre las APIs de Wompi, Rapyd (que adquirió PayU GPO), Mercado Pago y Kushki. Con ella, el comercio crea pagos con tarjeta o PSE, consulta su estado, lista los bancos PSE y valida las notificaciones asíncronas (webhooks).
  \item El módulo de navegador (\texttt{kit-pagos-colombia/browser}): obtiene el token de la tarjeta directamente de Wompi, Mercado Pago o Kushki desde la página del comprador, así que el número nunca llega al servidor del comercio. Con Rapyd, que cobra la tarjeta en su propia página, responde \texttt{UNSUPPORTED\_OPERATION}.
  \item La API de Simulación: un servicio Fastify desplegado en Render que tiene dos superficies: \texttt{/v1/sim} reproduce las APIs de las cuatro pasarelas, para probar sin credenciales reales, y \texttt{/v1/api} ofrece las operaciones del kit por REST. Desde el SDK, sin cabeceras especiales, alcanza los desenlaces de rechazo, pendiente, tiempo de espera y error, y emite los webhooks firmados de cada pasarela (RF-10 y RF-12 de \texttt{traceability-matrix.md}).
  \item La documentación centralizada de datos de prueba (\texttt{docs/testing-data/}): consolida las credenciales de sandbox, tarjetas, documentos, bancos, montos y flujos medidos de cada pasarela, con el nivel de evidencia de cada afirmación.
\end{itemize}

El framework contempla tres patrones de integración. En P1, el servidor del comercio usa el SDK. En P2, un cliente de navegador consume \texttt{/v1/api} sin instalar el SDK. En P3, el módulo de navegador unifica la tokenización de tarjeta; está implementado en parte, porque Rapyd no tokeniza en el navegador. Para demostrarlos y medirlos se construyeron cuatro proyectos prototípicos (\texttt{prototypes/}): A y B usan el SDK en el servidor (P1), C consume \texttt{/v1/api} desde el navegador (P2) y D integra las cuatro pasarelas a mano, como grupo de control. El proyecto se presenta además en una página web pública (\texttt{landing/}).

Los métodos de pago cubiertos son tarjeta y PSE. Quedan fuera el efectivo, Nequi, Daviplata, las suscripciones y los reembolsos. Tampoco se incluyen el manejo del número de tarjeta en el servidor del comercio ni otras pasarelas. También queda fuera actuar como procesador de pagos: Kit Pagos Colombia no autoriza, liquida ni procesa cobros por su cuenta. Arma cada petición, la envía a la pasarela (Wompi, Mercado Pago, Kushki o Rapyd) y normaliza la respuesta; el procesamiento lo hace siempre la pasarela. Las operaciones que el framework sí realiza, como crear pagos, consultar estados y verificar webhooks, forman parte del alcance. Quedan fuera, además, la integración con sistemas contables o ERP, la conciliación financiera, las operaciones administrativas sobre las cuentas de los proveedores y la responsabilidad sobre la continuidad de sus APIs. El proyecto tampoco incluye interfaces de pago para producción: las tiendas de los prototipos y el cliente del prototipo C son demostraciones contra los sandboxes. La protección de la página donde se escribe la tarjeta (su Content-Security-Policy) y el cálculo del monto son responsabilidad del comercio. La capa de seguridad propia de cada pasarela (autenticación con tokens, cifrado de datos, cumplimiento PCI-DSS) tampoco es modificada; el framework delega hacia cada proveedor sin interferir en sus mecanismos internos.

\subsubsection{Objetivos}

El objetivo general es desarrollar el Kit Pagos Colombia para simplificar la integración de múltiples pasarelas de pago en el contexto colombiano.

Los objetivos específicos son seis.

\begin{itemize}
  \item \textbf{Analizar} las características técnicas, mecanismos de integración y consideraciones de seguridad de las cuatro pasarelas seleccionadas.
  \item \textbf{Diseñar} la arquitectura del framework aplicando la arquitectura hexagonal y patrones GoF.
  \item \textbf{Desarrollar} un prototipo funcional del framework que implemente la interfaz unificada y los mecanismos de comunicación segura.
  \item \textbf{Implementar} la API de simulación.
  \item \textbf{Consolidar} los datos de prueba por pasarela para su documentación.
  \item \textbf{Evaluar} el framework mediante su aplicación en proyectos prototípicos usando las métricas WMC, CBO y RFC de Chidamber y Kemerer (1994), comparando los resultados frente a una integración directa con las APIs nativas.
\end{itemize}

\subsection{Supuestos y Restricciones}

\subsubsection{Supuestos}

El proyecto asume que las APIs de las cuatro pasarelas se mantienen estables en sus versiones actuales durante el período de desarrollo. Un cambio incompatible hacia atrás en alguna de ellas obligaría a rediseñar el adaptador correspondiente, lo que impactaría el cronograma. Para PayU, el supuesto ya no se cumplía cuando se escribió la versión 1.0 de este plan: Rapyd había completado la adquisición de PayU GPO el 14 de marzo de 2025, y el proyecto reemplazó a PayU por Rapyd.

Se asume que los entornos sandbox de las cuatro pasarelas están disponibles para consulta y prueba durante todas las fases del proyecto. Lo estuvieron, pero las mediciones mostraron que no permiten varios flujos de punta a punta. En Wompi, la URL del banco de PSE y el estado final del pago llegan en el mismo instante. En Mercado Pago, las credenciales de prueba reciben \texttt{401} en la Orders API, la única que cobra PSE. En Kushki, los cobros síncronos con tarjeta no quedan registrados para consulta. Esos flujos se prueban contra la API de Simulación.

Se asume que los cuatro integrantes del equipo dedicarán el tiempo acordado al proyecto durante el semestre, y que el director mantendrá disponibilidad para las reuniones de seguimiento establecidas. También se asume que es técnicamente viable construir una capa de abstracción suficientemente homogénea sobre las cuatro pasarelas para las operaciones fundamentales; esto está respaldado por trabajos previos como el de Manuaba e Hisyam (2022), aunque su validez específica para el conjunto colombiano se verificará en la fase de análisis.

\subsubsection{Restricciones}

El proyecto está acotado a un semestre académico, lo que limita la evaluación a proyectos prototípicos construidos por el propio equipo. El framework se distribuye bajo la licencia Apache 2.0 (\texttt{LICENSE} y campo \texttt{license} de \texttt{sdk/package.json}). Esta licencia permite usar, modificar y distribuir software libremente, incluyendo fines comerciales y creaciones propietarias, sin obligación de liberar el código fuente de las modificaciones, siempre que se incluyan los avisos de derechos de autor y cambios (Apache Software Foundation, 2004). La API de Simulación se despliega en Render, en el plan gratuito (\texttt{render.yaml}), lo que impone restricciones de disponibilidad y rendimiento que deben considerarse durante las pruebas. El SDK exige Node.js 20.19 o superior (campo \texttt{engines} de \texttt{sdk/package.json}). La elección del lenguaje de programación del SDK impone restricciones sobre las herramientas disponibles para calcular las métricas CK, aspecto que se aborda en la sección de Lenguajes y Herramientas.

\subsection{Entregables}

\begin{center}
\renewcommand{\arraystretch}{1.2}
\begin{longtable}{|P{0.227}|P{0.229}|P{0.519}|}
  \hline
  \textbf{Entregable} & \textbf{Estándares asociados} & \textbf{Justificación} \\
  \hline
  \endfirsthead
  \hline
  \textbf{Entregable} & \textbf{Estándares asociados} & \textbf{Justificación} \\
  \hline
  \endhead
  \caption{Entregables.}
  \label{tab:entregables}\\
  \endlastfoot
  SDK de servidor y módulo de navegador con interfaz unificada para Wompi, Rapyd, Mercado Pago y Kushki & Cockburn, A. (2024). \textit{Hexagonal Architecture Explained}; Gamma, E., Helm, R., Johnson, R., \& Vlissides, J. \textit{Design Patterns: Elements of Reusable Object-Oriented Software}; Hohpe, G. \& Woolf, B. (2003). \textit{Enterprise Integration Patterns} & La arquitectura hexagonal establece el estilo estructural del SDK, garantizando que el núcleo de lógica de pagos sea independiente de los proveedores. Los patrones GoF (Strategy, Facade y Factory) fundamentan las decisiones de diseño de la interfaz unificada y los adaptadores. Los patrones de integración empresarial de Hohpe y Woolf justifican el mecanismo de normalización de respuestas y el manejo de comunicaciones asíncronas entre el framework y las pasarelas. \\
  \hline
  API de Simulación desplegada en Render, con las rutas que imitan a las pasarelas (\texttt{/v1/sim}) y la API REST del kit (\texttt{/v1/api}) & Fielding, R. T. (2000). \textit{Architectural Styles and the Design of Network-based Software Architectures}; Ammann, P., \& Offutt, J. (2016). \textit{Introduction to Software Testing} & Los principios REST de Fielding fundamentan el diseño de la interfaz de la API de simulación, garantizando una interfaz coherente, sin estado y desacoplada de cualquier cliente. Los criterios de Ammann y Offutt justifican la cobertura de escenarios que debe exponer la API (aprobación, rechazo, pendiente, tiempo de espera, error y notificaciones asíncronas) y el entorno controlado equivalente a un sandbox. \\
  \hline
  Documentación centralizada de datos de prueba por pasarela & Ammann, P., \& Offutt, J. (2016). \textit{Introduction to Software Testing} & La noción de entorno sandbox como espacio de verificación aislado, definida por Ammann y Offutt, justifica la necesidad de consolidar los datos de prueba de cada proveedor en un formato accesible y estructurado, reduciendo la dependencia de fuentes dispersas durante el desarrollo. \\
  \hline
  Informe comparativo de la evaluación: seis variables medidas sobre los prototipos B (con el framework) y D (integración directa), entre ellas WMC, CBO y RFC & Chidamber, S. R. \& Kemerer, C. F. (1994). \textit{A metrics suite for object oriented design}. IEEE Transactions on Software Engineering, 20(6), 476–493 & Las métricas CK (Chidamber y Kemerer) proveen una herramienta cuantitativa y reproducible para comparar la complejidad, el acoplamiento y la superficie de interacción del código de integración con y sin el framework. Su uso garantiza que la evaluación del framework no dependa exclusivamente de observaciones subjetivas. \\
  \hline
  Documento final del trabajo de grado y repositorio público con documentación de uso & Peffers, K., Tuunanen, T., Rothenberger, M., \& Chatterjee, S. (2007). \textit{A design science research methodology for information systems research}. Journal of Management Information Systems, 24, 45–77 & DSR exige que los artefactos producidos y las lecciones aprendidas sean comunicados de forma rigurosa a las audiencias relevantes. El documento de grado y el repositorio público materializan esta exigencia, presentando el proceso de construcción, los resultados de la evaluación y las condiciones bajo las cuales el framework resulta útil. \\
  \hline
\end{longtable}
\end{center}

\subsection{Evolución del Plan}

Este documento no es estático. A medida que el proyecto avance y se produzcan ajustes en el alcance, el cronograma, los riesgos o las decisiones técnicas, el SPMP deberá actualizarse para reflejar el estado real del proyecto. Cada entrega parcial del trabajo de grado es una oportunidad obligatoria de revisión, especialmente de la carta Gantt descrita en la sección de \textbf{Calendarización}, que debe estar actualizada en cada entrega.

Cualquier integrante puede proponer un cambio al plan. La propuesta se discute en la reunión semanal del equipo y, si afecta el alcance o los criterios de evaluación definidos con el director, se comunica antes de implementarse. Los cambios aprobados se registran en el historial de cambios del documento con fecha, autor y descripción breve. El SPMP se versiona en el repositorio del proyecto bajo control de cambios en GitHub.

Los cambios que no requieren aprobación del director son: ajustes al cronograma que no afecten los hitos principales, actualización de la tabla de riesgos, correcciones menores de redacción y actualización de métricas de progreso. Los que sí requieren alineación son:

\begin{itemize}
  \item Modificaciones al alcance del proyecto
  \item Cambios en los entregables o sus criterios de aceptación
  \item Cambios en la metodología de evaluación.
\end{itemize}

Se anticipa al menos una actualización por cada hito completado, con una revisión mayor al finalizar la fase de desarrollo para reflejar los resultados de las pruebas internas y los ajustes que hayan surgido durante la implementación de los adaptadores.
